\section{Introducción}
    Los sistemas lineales invariantes en el tiempo (LIT) son fundamentales en la teoría de sistemas y el análisis de señales debido a que cumplen que la respuesta a una combinación lineal de señales es una combinación lineal de las respuestas individuales, y además a su invarianza temporal, es decir, sus propiedades no cambian con el tiempo \cite{sistemasLIT}. Analizar la respuesta de estos sistemas en función de la frecuencia de la señal de entrada posee una gran ventaja, si se estimulan con una señal sinusoidal, la respuesta será una señal de la misma forma pero con diferente amplitud y fase, siendo la manera en la que se modifica la señal dependiente de la frecuencia y las características del sistema. De esta manera, resulta posible caracterizar el comportamiento del sistema según las diferencias entre la señal de entrada y su respuesta. En particular, los circuitos RC y RLC cumplen estas propiedades.

    Al alimentar un circuito RC con una señal de corriente alterna sinusoidal \(V_e\) de frecuencia angular \(\omega\), la respuesta en tensión del circuito \(V_s\) será de la forma
    \begin{equation}
        V_s = H(\omega) V_e
        \label{cociente}
    \end{equation}
    donde \(H(\omega)\) es la función de transferencia del sistema, que depende de la frecuencia de la señal de entrada y las características del circuito. Típicamente, a esta función se la modela como un número complejo, permitiendo así representar tanto la atenuación como el desfasaje de la señal. En el caso del circuito RC, este funciona como un filtro de frecuencias bajas visto desde la resistencia (filtro PA) y como un filtro de frecuencias altas visto desde el capacitor (filtro PB), lo que implica que la expresión de la función de transferencia será diferente dependiendo de donde se mida la señal de salida. Al medir la tensión de salida en la resistencia, el módulo de \(H(\omega)\) tendrá la forma \cite{oppenheim1997signals}
    \begin{equation}
        |H(\omega)| = \frac{\omega \tau}{\sqrt{1 + (\omega \tau)^2}} 
     \label{PA}
    \end{equation}
    y su fase 
    \begin{equation}
        \phi = \arctan\left(\frac{1}{\omega \tau}\right)
        \label{eq:3}
    \end{equation}
    donde se definió \(\tau = RC\) como la constante de tiempo del circuito \cite{kip1969}. En cambio, si se mide la tensión de salida en el capacitor, el módulo de la función de transferencia tendrá la forma 
    \begin{equation}
        |H(\omega)| = \frac{1}{\sqrt{1 + (\omega \tau)^2}}
        \label{PB}
    \end{equation}
    y su fase
    \begin{equation}
        \phi = -\arctan(\omega \tau)
        \label{eq:4}
    \end{equation}
    En ambos casos, se puede observar que el módulo de la función de transferencia refleja la característica de filtro de cada componente. En el caso del capacitor, la función de transferencia crece en módulo a medida que la frecuencia de la señal de entrada disminuye mientras que, en el caso de la resistencia, el módulo crece a medida que la frecuencia de la señal de entrada aumenta. Además, resulta de interés definir una frecuencia de corte como aquella a la cual la potencia cae a la mitad de su valor máximo, lo que implica que el módulo de la función de transferencia es \(H(\omega) = \frac{1}{\sqrt{2}}\). Al aplicar esta condición, se obtiene que la frecuencia de corte es \(\omega_0 = \frac{1}{\tau}\) en ambos casos.


De igual manera, el circuito RLC presenta una función de transferencia que caracteriza la respuesta ante señales de diferente frecuencia. En este caso, el circuito actúa como un filtro selectivo de frecuencias cercanas a la frecuencia de resonancia del sistema \(\omega_r\). De esta manera, el módulo de la función de transferencia del circuito RLC se expresa como \cite{oppenheim1997signals}
    \begin{equation}    
        |H(\omega)| = \frac{\omega \tau}{\sqrt{(1 - LC\omega^2)^2 + (\omega \tau)^2}}
        \label{eq:moduloRLC}
    \end{equation}  
    y su fase
    \begin{equation}
        \phi = \arctan\left(\frac{ LC\omega^2 - 1}{\omega \tau}\right)
        \label{eq:desfasajeRLC}
    \end{equation}
    donde al encontrar el máximo del módulo de la función de transferencia, se obtiene que la frecuencia de resonancia es \(\omega_r = \frac{1}{\sqrt{LC}}\). Idealmente, en esta frecuencia el inductor y el capacitor dejan de aportar al circuito y se comportan como cables ideales, dejando a la resistencia como el único elemento del circuito \cite{TiplerMosca}.

    Con el objetivo de determinar las constantes características de ambos circuitos y analizar su respuesta en frecuencia, se construyeron diagramas de Bode y Nyquist a partir de medidas de tensión de entrada y salida, junto al desfasaje entre ambas señales.